% Nouveau EGA V (cote 25): archive cote 25, site page 64 (PDF page 65).
% typescript. Draft compared: none.
% AI-transcribed by Claude Opus 5.5 (claude-opus-5-5, adaptive thinking, effort high), Message Batches API, batch ega25b (msgbatch_01A7n5YuXubUPNK1zsepi8RG), 2026-11-13; not yet checked by eye.
\scanpage{25/64}
\mspage{\unsure{57}}

point rationnel sur k.

Un point de $\underline{\mathrm{Grass}}_n(\underline{E})$ à valeurs dans S, i.e. un Module quotient
loc libre de rang n \ $\underline{G}$ de $\underline{E}$, définit canoniquement un sous-préschéma
de $P(\underline{E})$ , savoir $P(\underline{G})$ . Un tel sous-préschéma\inserted{(mais sans imposer le rang de $\underline{G}$)} est appelé \uline{sous-variété}
% NOTE: the last letter of "sous-variété" is cut by the right edge of the scan
\uline{linéaire} de $P(\underline{E})$ (relativement à S si une confusion e st possible).
C'est donc un fibré projectif,de dimension relative n-1 si $n\geqslant 1$ ,
(et vide si n=0). On vérifie immédiatement que la section de $\underline{\mathrm{Grass}}_n$
$(\underline{E})$ i.e. G \ est connu quand on connait la sous-variété linéaire
correspondante de $P(\underline{E})$. D\unsure{e} cette façon, \ la grassmanienne peut
s'interpréter comme représentant le foncteur "sous-variétés linéaires
de dimension relative n-1 \xout{Xx=xresp xxxvide=xxx} de $P_{S'}$ " pour S'
variable\xout{x}\inserted[typed]{$(n\geqslant 1)$}. Il est d'ailleurs possible de donner une caractérisation
intrinsèque de ce dernier foncteur i.e. de la notion de sous-variété
linéaire \ de dimension relative m : ce sont les sous-préschémas fermés
de P , lisses sur S et partout de dimension relative m, et qui sont
de "degré projectif 1" \ en tout $s\in S$; \ cette caractérisation sera
donnée dans un Chapitre ultérieur, et nous n'en avons nul besoin ici.

Supposant de nouveau $\underline{E}$ \ loc libre\inserted[typed]{de rang r}, soit $\underline{E}$ \ son dual. Alors
par polarité nous trouvons un isomorphisme canonique
\[ \underline{\mathrm{Grass}}_n(E) \simeq \underline{\mathrm{Grass}}_{r-n}(\check{\underline{E}}) \]
faisant correspondre \ au \ quotient \ $\underline{G}$ de $\underline{E}$ \ le quotient \ $\check{\underline{E}}/\check{\underline{G}}$ \ de $\check{\underline{E}}$ .
Du point de vue variétés linéaires, à \ la variété linéaire de dim rel
m \ L dans P, correspond la \ variété linéaire\inserted[typed]{duale $L^{o}$} de dim relative
\xout{(x=1)xmx} \ (r-1)-1 -m \ \ de $\check{P}$,\inserted{i.e. de codim relative m+1 dans $\check{P}$} (NB \ r-1 \ est ici la dim relative
commune de \ P et $\check{P}$ sur S), qu'on peut visualiser géométriquement ainsi,
Faisant d'abord \ n=r-1, on trouve \ un isomorphisme
\[ \underline{P}(\check{\underline{E}}) \simeq \underline{\mathrm{Grass}}_{r-1}(\underline{E}) \]

\mspage{64}
